\documentclass[11pt]{article}
\usepackage[a4paper,margin=29mm]{geometry}
\usepackage[T1]{fontenc}
\usepackage{lmodern,microtype,amsmath,amssymb,amsthm,mathtools}
\usepackage{booktabs,array,enumitem}
\usepackage{xcolor}
\usepackage[colorlinks=true,linkcolor=blue!45!black,citecolor=blue!45!black,urlcolor=blue!45!black]{hyperref}
\hypersetup{pdftitle={Temporal Sampling and Local Nonidentifiability in Coalescent Demography},pdfauthor={KEIO}}
\newtheorem{theorem}{Theorem}[section]
\newtheorem{lemma}[theorem]{Lemma}
\newtheorem{proposition}[theorem]{Proposition}
\newtheorem{corollary}[theorem]{Corollary}
\theoremstyle{definition}
\newtheorem{definition}[theorem]{Definition}
\newtheorem{example}[theorem]{Example}
\theoremstyle{remark}
\newtheorem{remark}[theorem]{Remark}
\newcommand{\R}{\mathbb R}
\newcommand{\C}{\mathbb C}
\newcommand{\N}{\mathbb N}
\newcommand{\E}{\mathbb E}
\newcommand{\dd}{\,d}
\newcommand{\supp}{\operatorname{supp}}
\newcommand{\rank}{\operatorname{rank}}
\newcommand{\Lip}{\operatorname{Lip}}
\newcommand{\TV}{\operatorname{TV}}
\newcommand{\W}{W_1}
\DeclareMathOperator{\diag}{diag}
\setlist[enumerate]{leftmargin=*,itemsep=3pt}
\setlist[itemize]{leftmargin=*,itemsep=3pt}
\setlength{\emergencystretch}{2em}
\numberwithin{equation}{section}
\title{Temporal Sampling and Local Nonidentifiability\\in Coalescent Demography}
\author{KEIO}
\date{20 October 2026}
\begin{document}
\maketitle
\begin{abstract}
We study which temporal sampling designs identify a nonparametric
population-size history from absolute expected site-frequency spectra
under the variable-size Kingman coalescent. On positive, uniformly bounded
smooth histories, the collection of all expected joint serial spectra,
with arbitrary finite sample sizes at dates in a set $S$, is injective
if and only if $S$ is dense. Every sampling gap supports an
infinite-dimensional family of smooth calendar-local perturbations that
preserve all these expectations and the ancestral partition skeleton at
the allowed dates. The perturbations can be arbitrarily small in any
fixed finite-order smooth norm and preserve strict finite-order shape
constraints. We construct them by a compact Gevrey realization of the
classical Kingman spectral annihilator. We then give two ways to remove
the obstruction. Endpoint-resolved occupation rewards expose the
consecutive differences of the quadratic coalescent spectrum, recovering
a complete Hausdorff moment sequence; finitely many such correlations
have a sharp inverse-linear Wasserstein resolution. Complete monotonicity
of the inverse clock also restores identifiability from contemporaneous
spectra through positive-real interpolation. Dense serial sampling further
identifies an unknown mutation intensity together with demography.
The statements concern the specified expectation and reward observations,
not equality of full continuous-time genealogy or linked-genome laws.
\end{abstract}
\noindent\textbf{Keywords.} Coalescent theory; demographic identifiability;
serial sampling; site-frequency spectrum; Gevrey functions; inverse moments;
completely monotone functions.

\section{Introduction}

The amount of genetic data and the temporal placement of the sampled
individuals are different sources of information. Increasing the number
of individuals at fixed dates probes a larger collection of coalescent
rates. Moving the dates changes which portions of the demographic clock
are directly accessible. This paper separates these mechanisms in an
idealized exact-expectation model.

The classical Kingman coalescent has exit rates
$\lambda_k=\binom{k}{2}$ when $k$ ancestral lineages are present
\cite{Kingman}. A demographic history changes the speed of this process.
In coalescent time, the unknown quantity is the derivative of the inverse
clock. Expected branch lengths are linear functionals of that derivative,
with kernels built from $e^{-\lambda_k u}$. The quadratic spacing of the
rates permits nonzero smooth perturbations annihilating every such
functional. Serial sampling adds lineages at prescribed calendar times;
it does not, within a sampling gap, change this spectral set.

Our first result makes that observation local and uniform over every
finite serial design. If an open interval contains no allowed sampling
date, the history can be changed only inside that interval while
preserving all expected serial spectra, however large the sample sizes.
The perturbation has zero integral in coalescent time, so it preserves
both calendar endpoints of the interval. This last condition also
preserves the evolution of a single ancestral lineage before additional
ancient samples are inserted. Conversely, dense dates determine the mean
residual coalescence time of two contemporaneous lineages as a continuous
function. Its first derivative recovers the population size.

Our second result changes the observable rather than the temporal design.
For a fixed interval, we multiply the time spent in the initial
$n$-lineage state by the indicator that exactly one merger has occurred
at the endpoint. The expectation of this product exposes
$\lambda_n-\lambda_{n-1}=n-1$, rather than only $\lambda_n$.
Consequently these scalar correlations determine all ordinary moments of
an explicit probability measure and identify the clock. They constitute
additional genealogical information; they are not ordinary expected-SFS
coordinates. The finite-correlation ambiguity is quantified by matching
upper and lower Wasserstein bounds of order $1/m$.

A third result identifies a structural restriction under which the sparse
quadratic observations already suffice. Completely monotone inverse
clocks give Stieltjes transforms. A square-root change of variable turns
these into positive-real functions on a half-plane, for which the
quadratic sampling sequence becomes a uniqueness set. This supplies a
precise contrast between preserving finitely many strict derivative
inequalities and imposing complete monotonicity. We also show how
heterochronous pairs calibrate a constant mutation rate, and how dense
dates jointly identify a time-dependent mutation intensity.

\subsection{Relation to prior results}

Myers, Fefferman, and Patterson \cite{MFP} established the basic
nonidentifiability from contemporaneous spectra. Their equation (21)
already contains the factor
$\cos((\pi/2)\sqrt{1+8z})$ used below, and their construction includes
smooth perturbations in coalescent time. We do not claim this factor or
the original nonidentifiability phenomenon as new. Here it is realized
as an infinite-order differential operator on a compact Gevrey bump,
which permits support inside any prescribed sampling gap. The
calendar-local serial-design theorem and the endpoint-conditioned
restoration result are the specific extensions developed in this paper.

Bhaskar and Song \cite{BS} prove identifiability for demographic families
with controlled sign-change complexity, including fixed-complexity
piecewise models. Such control concerns differences of rescaled histories;
it is stronger than requiring each history to be monotone. Our
perturbations retain unbounded sign-change complexity in their
differences. Spence, Kamm, and Song \cite{SKS} give spectral formulas and
identifiability results for general coalescents. We use the standard
Kingman process and prove the partition-generator statements required
for serial rewards directly.

The polynomial approximation and positive-real-function arguments use
classical tools. Their purpose here is to identify which additional
observations or shape restrictions remove the particular coalescent
obstruction. No result below equates expected spectra with full
linked-sequence likelihoods. The paper establishes the stated mathematical
claims; it does not assert a resolution of a named open conjecture.

\subsection{Outline}

Section~\ref{sec:model} states the temporal-design theorem.
Sections~\ref{sec:kingman}--\ref{sec:dense} prove its two directions.
Section~\ref{sec:consequences} records local shape and statistical
consequences. The subsequent sections establish recovery from
endpoint-resolved rewards, complete monotonicity, and mutation-rate
calibration. Appendix~\ref{sec:sparse} places the compact annihilator in
a general spectral framework and proves the complementary obstruction
for reciprocally divergent spectral sequences.



\section{Observation model and the temporal-design criterion}
\label{sec:model}

Time is measured backwards from the present. Let $\mathcal H$ be the
class of functions $\eta\in C^\infty([0,\infty))$ such that
$0<\inf\eta\le\sup\eta<\infty$. For a given history $\eta$, any pair
of currently existing ancestral lineages merges at rate $1/\eta(t)$.
At each prescribed sampling date, the corresponding labeled ancestral
lineages are inserted. The sampling schedule is deterministic and is
not itself a random observation informative about $\eta$.

For a finite labeled sampling schedule $\mathcal D$, let
$\mathcal B_{\mathcal D}(\eta)$ denote the vector of expected lengths
of branches subtending each nonempty proper subset of the sampled labels.
The expected joint serial site-frequency spectrum is obtained by summing
these coordinates according to the numbers of descendants in each
sampling cohort and multiplying by a known, constant mutation rate.
We retain absolute expected branch lengths, not only a normalized
spectrum conditional on a mutation.

For $S\subseteq[0,\infty)$, define the observation map
\[
 \mathcal O_S(\eta)=
 \bigl(\mathcal B_{\mathcal D}(\eta):
        \mathcal D\text{ is a finite sampling schedule with dates in }S\bigr).
\]
Sample sizes at the allowed dates are unrestricted.

\begin{theorem}[Exact criterion for the time design]\label{phy:design}
The map $\mathcal O_S$ is injective on $\mathcal H$ if and only if
$S$ is dense in $[0,\infty)$.
The same criterion holds if $\mathcal B_{\mathcal D}$ is replaced by
the absolute expected joint serial site-frequency spectrum.

More precisely, if a nonempty open interval $I\subset(0,\infty)$
contains no point of $S$, then every $\eta_0\in\mathcal H$ belongs
to an infinite-dimensional family of distinct histories with identical
$\mathcal O_S$, all equal to $\eta_0$ outside $I$. For every fixed
integer $r\ge0$, members of this family may be chosen arbitrarily
close to $\eta_0$ in $C^r$. Any finite collection of strict differential
inequalities satisfied by $\eta_0$ with a positive margin on a compact
subinterval of $I$ can be preserved.
\end{theorem}

The negative assertion does not concern equality of full genealogy
distributions, full linked-sequence likelihoods, or distributions of
the random spectrum at a finite locus. It concerns exact expected
branch rewards, simultaneously for all finite schedules on $S$.
In a statistical model postulating independent Poisson counts with
these expected spectra as intensities, it also gives equality of the
entire observation law for every exposure.

\section{The Kingman spectral annihilator}
\label{sec:kingman}

Write $\lambda_k=k(k-1)/2$ for $k\ge1$. In particular $\lambda_1=0$.
The spectral factor below is the one in \cite[equation (21)]{MFP}.
We realize it on a compact Gevrey bump so that the resulting perturbation
fits inside an arbitrary prescribed interval.

\begin{lemma}[A smooth local nullspace]\label{phy:nullspace}
For every bounded open interval $J\subset\mathbb R$, the space
\[
 \mathcal N_J=\left\{h\in C_c^\infty(J;\mathbb R):
      \int_{\mathbb R}h(u)e^{-\lambda_k u}\,du=0
                        \text{ for every }k\ge1\right\}
\]
is infinite-dimensional.
\end{lemma}
\begin{proof}
The function
\[
 P(z)=\cos\left(\frac\pi2\sqrt{1+8z}\right)
     =\sum_{j=0}^\infty
       \frac{(-1)^j(\pi^2/4)^j(1+8z)^j}{(2j)!}
\]
is entire: the expression with a square root is defined by the displayed
power series and is independent of its branch. It vanishes at every
$\lambda_k$, since $\sqrt{1+8\lambda_k}=2k-1$.
There are constants $C_0,C_1$ such that
$|P(z)|\le C_0e^{C_1\sqrt{|z|}}$.
Writing $P(z)=\sum_{j\ge0}a_jz^j$, Cauchy's estimate on the circle of
radius $(2j/C_1)^2$, for $j\ge1$, gives
\[
 |a_j|\le C_0\left(\frac{eC_1}{2j}\right)^{2j}
          \le C_0 B^j/(j!)^2
\]
for a fixed $B$.

Choose a nonnegative nonzero compactly supported Gevrey function
$\kappa$ in $J$ of order $\sigma$ with $1<\sigma<2$; thus
$\|\kappa^{(j)}\|_\infty\le C A^j(j!)^\sigma$.
For completeness, order $3/2$ is sufficient and elementary: the function
$f(x)=e^{-1/x^2}$ for $x>0$, extended by zero for $x\le0$, has
local derivative bounds of that order. Cauchy's estimate on a circle
of radius $\delta x$ bounds its $j$th derivative by
$j!(\delta x)^{-j}e^{-c/x^2}$; maximizing over $x$ gives a bound
$C^j(j!)^{3/2}$. A product $f(u-a)f(b-u)$, where
$[a,b]\subset J$, provides the required bump.

The series
\[
                     h=P(D)\kappa=\sum_{j\ge0}a_j\kappa^{(j)}
\]
converges in every $C^r$ norm. Indeed its $r$th derivative is bounded
termwise by $C_rB_r^j/(j!)^{2-\sigma}$, a summable sequence.
All terms have support in the same compact subset of $J$, hence
$h\in C_c^\infty(J)$. Integration by parts and locally uniform
convergence give, for every complex $z$,
\[
 \int h(u)e^{-zu}\,du
       =P(z)\int\kappa(u)e^{-zu}\,du.
\]
This proves the vanishing conditions, including $\int h=0$.
It also proves $h\ne0$: at $z=-1$, both
$P(-1)=\cosh(\pi\sqrt7/2)$ and the transform of $\kappa$ are positive.

Every derivative $h^{(j)}$ belongs to $\mathcal N_J$. These derivatives
are linearly independent. Otherwise their Laplace transforms would
give $Q(z)\widehat h(z)=0$ for a nonzero polynomial $Q$, which is
impossible for the nonzero entire function $\widehat h$.
\end{proof}

Appendix~\ref{sec:sparse} proves the corresponding local construction for
superlinear spectra with bounded Jordan multiplicities, and a uniqueness
obstruction for reciprocally divergent spectra.

\subsection{Why all labeled branch rewards have the same spectrum}

\begin{lemma}[No hidden Jordan factors]\label{phy:semisimple}
Fix a finite set of currently sampled labels and let $Q$ be the
generator of its partition-valued Kingman coalescent in coalescent
time. Then
\[
                   e^{uQ}=\sum_{k=1}^m e^{-\lambda_k u}\Pi_k
\]
for constant matrices $\Pi_k$, some of which may be zero.
Consequently every coordinate of $v e^{uQ}r$, for fixed vectors
$v,r$, is a linear combination of these exponentials.
\end{lemma}
\begin{proof}
Order partitions by their number of blocks. The diagonal block at
level $k$ is $-\lambda_k I$, and the only off-diagonal transitions
are from level $k$ to level $k-1$. Distinct levels have distinct
$\lambda_k$. A triangular block matrix whose diagonal blocks are
scalar matrices with pairwise disjoint spectra is similar to its
block-diagonal part: successively eliminate each off-diagonal block
by the Sylvester equation
$(-\lambda_k)X-X(-\lambda_j)=B$, whose solution is
$X=B/(\lambda_j-\lambda_k)$. Thus $Q$ is diagonalizable and has the
stated spectral expansion. Repeated eigenvalues across partitions
at the same level introduce no Jordan blocks.
\end{proof}

\section{Calendar-local perturbations and serial rewards}
\label{sec:gap}

For a baseline history $\eta_0$, define its coalescent clock and its
inverse by
\[
 C_0(t)=\int_0^t\frac{dv}{\eta_0(v)},\qquad
 T_0=C_0^{-1},\qquad f_0(u)=T_0'(u)=\eta_0(T_0(u)).
\]
Both clocks are global smooth increasing diffeomorphisms, by the
uniform positive upper and lower bounds on $\eta_0$.
Choose $0<a<b<\infty$ with $[a,b]\subset I$, and take
$h\in\mathcal N_{(C_0(a),C_0(b))}$. Set
\[
 H(u)=\int_0^u h(v)\,dv,\qquad
 f_\varepsilon=f_0+\varepsilon h,\qquad
 T_\varepsilon=T_0+\varepsilon H,
 \qquad \eta_\varepsilon(T_\varepsilon(u))=f_\varepsilon(u).
\]
For sufficiently small $|\varepsilon|$, $f_\varepsilon$ stays uniformly
positive, so these formulas define a history in $\mathcal H$.
Because $\int h=0$, both $h$ and $H$ vanish outside the same compact
coalescent-time interval. Hence $T_\varepsilon=T_0$ and
$\eta_\varepsilon=\eta_0$ outside $(a,b)$ in calendar time as well.
The map $h\mapsto\eta_h$ is injective: a calendar history uniquely
determines its inverse clock derivative $f$.

Fix now an arbitrary finite sampling schedule on $S$. Its dates
avoid $(a,b)$ and their coalescent-clock locations are unchanged.
If no lineage has yet been inserted, the affected interval contributes
zero reward. Otherwise, on every sampling epoch the ancestral partition evolves with the
same homogeneous generator $Q$. In particular, the transition
matrix across $(a,b)$ depends only on
$C_0(b)-C_0(a)$, which is unchanged. The law of the partition at
every calendar time outside $(a,b)$ is therefore unchanged, by
induction over the sampling dates. More strongly, all finite-dimensional
laws of its partition skeleton at those times agree, since every
transition across the modified interval has the same coalescent duration.

Within the affected interval, let $v$ be the entering partition law
and let $r$ be any state-dependent branch reward. Writing
$A=C_0(a)$ and $B=C_0(b)$, its expected calendar-time integral is
\[
 \int_A^B v e^{(u-A)Q}r\,f_\varepsilon(u)\,du.
\]
By Lemma~\ref{phy:semisimple}, the difference from its baseline value
is a linear combination of
\[
             \varepsilon\int_A^B e^{-\lambda_k u}h(u)\,du=0.
\]
This includes the level-one reward through the condition at
$\lambda_1=0$, which is important when further ancient samples
remain to be inserted. Taking $r$ to count branches with each
specified descendant-label set proves equality of all expected
labeled branch lengths and hence all expected joint serial spectra.
Future samples are not descendants of a branch before their insertion;
therefore these rewards depend only on the current labeled partition,
not on a future endpoint condition. No fixed upper bound on the sample
size was used. After the oldest insertion, the full-label branch has zero
reward, so the integrals are finite by boundedness of the history.

The finite-dimensional slices generated by
$h,h',\ldots,h^{(d-1)}$ yield an injective $d$-parameter family for
every $d$. Smooth dependence of inverse diffeomorphisms on compact
sets implies $\eta_h\to\eta_0$ in any prescribed $C^r$ norm as
$h\to0$ in the corresponding $C^r$ norm. Since the perturbation is
supported in a fixed compact subinterval of $I$, strict differential
inequalities with a positive margin there persist. For example,
if $\eta_0'>0$ throughout $[a,b]$, both distinct perturbed histories
can be chosen strictly increasing there and unchanged elsewhere.
This completes the negative assertion of Theorem~\ref{phy:design}.

\section{Dense dates recover the history from pairs}
\label{sec:dense}

Two contemporaneous lineages sampled at calendar time $t$ have mean
residual coalescence time
\[
 m_\eta(t)=\int_t^\infty
       \exp\left(-\int_t^v\frac{dw}{\eta(w)}\right)dv.
\]
The uniform bounds on $\eta$ make this finite and justify
differentiation under the integral. Directly,
\[
                 m_\eta'(t)=\frac{m_\eta(t)}{\eta(t)}-1,
 \qquad
                 \eta(t)=\frac{m_\eta(t)}{1+m_\eta'(t)}.
\]
The absolute expected pairwise spectrum determines $2m_\eta(t)$,
up to the known mutation-rate factor. Equality of observations on
a dense set $S$ therefore gives equality of the continuous functions
$m_\eta$ on all of $[0,\infty)$, and then equality of their derivatives
and of $\eta$. This proves the positive assertion.



\section{Local shape constraints and information boundaries}
\label{sec:consequences}

\begin{corollary}[Strict shape restrictions do not close a sampling gap]
\label{cor:shape}
Let $I$ be a sampling gap. Suppose $\eta_0\in\mathcal H$ satisfies
a finite collection of inequalities
\[
 F_j\bigl(t,\eta_0(t),\eta_0'(t),\ldots,\eta_0^{(r)}(t)\bigr)>0,
 \qquad 1\le j\le q,
\]
on a compact interval $J\subset I$ with nonempty interior, where the
$F_j$ are continuous. There are distinct histories arbitrarily close
to $\eta_0$ in $C^r$, equal to it outside $J$, with the same
$\mathcal O_S$ and satisfying the same strict inequalities on $J$.
\end{corollary}
\begin{proof}
On a slightly smaller compact interval the inequalities have a common
positive margin. Apply the gap construction with support in its interior
and choose the perturbation small in $C^r$. Uniform continuity preserves
the inequalities there; outside the support they are unchanged.
\end{proof}

For example, $\eta_0(t)=2+e^{-t}$ belongs to $\mathcal H$ and is strictly
decreasing and strictly convex. On every compact sampling gap, distinct
histories can be chosen with these same properties, and with any fixed
finite collection of the alternating derivative signs of this baseline.
This is a mathematical statement about finite-order shape restrictions.
It does not imply finite sign-change complexity of differences, nor
preservation of all derivative inequalities simultaneously.

\begin{proposition}[Finite linear calibration]
\label{prop:calibration}
Let $\ell_1,\ldots,\ell_q$ be any linear functionals defined on
$C_c^\infty(J)$, where $J$ is a coalescent-time interval inside a
sampling gap. There are arbitrarily large finite-dimensional spaces
of perturbations that preserve all expected serial spectra and also
satisfy $\ell_i(h)=0$ for $1\le i\le q$.
\end{proposition}
\begin{proof}
The nullspace in Lemma~\ref{phy:nullspace} is infinite-dimensional.
Intersecting it with the kernels of $q$ linear functionals reduces its
algebraic dimension by at most $q$. Any finite-dimensional subspace of
the intersection can be scaled to preserve positivity of the inverse
clock. This statement concerns linear calibration of that clock;
it makes no assertion about arbitrary nonlinear auxiliary measurements.
\end{proof}

\begin{example}[Infinitely many dates need not be enough]
The allowed dates $S=\{0,1,2,\ldots\}$ leave a nonidentifiable
perturbation inside $(1/3,2/3)$, even if all finite serial schedules
and all sample sizes are available. The set
$\{0\}\cup\{1/n:n\ge1\}$ is also insufficient. By contrast any
countable dense subset of $[0,\infty)$ is sufficient in the
exact-expectation model. Density is an ideal information condition;
it is not a finite-data sampling prescription.
\end{example}

\begin{proposition}[A statistical consequence under a specified count model]
\label{prop:poisson}
Suppose an experiment, after choosing a finite schedule on $S$ and
an exposure, returns independent Poisson counts whose intensities
are that exposure times the coordinates of its expected joint serial
spectrum. Two histories in the same gap family have identical laws
for every finite sequence of these experiments, including designs
chosen from previous observations. For any metric $d$ on histories,
the maximum expected loss of any estimator at two such histories
$\eta_0,\eta_1$ is at least $d(\eta_0,\eta_1)/2$.
\end{proposition}
\begin{proof}
Every possible next design has the same vector of intensities under the
two histories. Induction on the experiments gives the same joint law,
also for randomized choices based on the observed past. Under that common
law, the triangle inequality gives
\[
 \E d(\widehat\eta,\eta_0)+\E d(\widehat\eta,\eta_1)
 \ge d(\eta_0,\eta_1).
\]
At least one of the two risks is at least half the right-hand side.
\end{proof}

The Poisson model in Proposition~\ref{prop:poisson} is an explicitly
declared observation model. The random SFS at a linked locus can contain
covariance information beyond its mean. Equality of expected spectra
does not establish equality of that law. Likewise, the preserved
ancestral partition skeleton at allowed dates does not specify the
continuous-time merger locations inside a gap.



\section{Restoring a sampling-gap clock by endpoint--occupation correlations}
\label{restore:section}

The preceding nonidentifiability theorem concerns unconditioned expected
branch rewards.  We now specify a different observation that determines
the clock inside a single gap.  The new data are not ordinary expected
site-frequency-spectrum coordinates: an occupation time is multiplied by
an indicator of the terminal block count.  This small change replaces the
sparse coalescence-rate spectrum by all its adjacent differences.

Fix a calendar interval $[a,b]$ containing no sampling insertion, and write
$T=b-a>0$.  Relative coalescent time on this interval is
\[
 u(t)=\int_a^t\frac{dv}{\eta(v)},\qquad L=u(b),
 \qquad f(u)=\frac{dt}{du}>0,\qquad \int_0^L f(u)\,du=T.
\]
We assume $f\in C^\infty([0,L])$ and $f>0$ everywhere.  For an experiment
started with $n$ lineages at $a$, let $K^{(n)}_u$ be its block count in
coalescent time.  It is the pure-death chain with rate
$\lambda_k=\binom{k}{2}$ from $k$ to $k-1$.

Define the endpoint--occupation correlation
\begin{equation}\label{restore:Rdef}
 R_n(f)=\mathbb E\left[
     \left(\int_0^L f(u)\mathbf 1_{\{K^{(n)}_u=n\}}\,du\right)
                  \mathbf 1_{\{K^{(n)}_L=n-1\}}\right],\qquad n\ge2.
\end{equation}
The integral inside the expectation is calendar time spent at the initial
block count.  Expectations refer to the chain started in state $n$.

\begin{theorem}[Identification from adjacent endpoint correlations]
\label{restore:identify}
The no-merger probability for two initial lineages, together with $T$ and
the correlations $R_n(f)$ for all $n\ge2$, uniquely determines $L$, $f$,
and the calendar clock on $[a,b]$.

More explicitly, put $F(s)=\int_0^L f(u)e^{-su}\,du$.  The no-merger
probability is $e^{-L}$, and
\begin{equation}\label{restore:Rformula}
 R_n(f)=\frac{\lambda_n}{n-1}
       \left[e^{-\lambda_{n-1}L}F(n-1)-e^{-\lambda_nL}T\right].
\end{equation}
Consequently, for every integer $j\ge1$,
\begin{equation}\label{restore:Fformula}
 F(j)=\frac{j}{\lambda_{j+1}}e^{\lambda_jL}R_{j+1}(f)
                         +e^{-jL}T.
\end{equation}
\end{theorem}
\begin{proof}
For $v\ge0$, the probability of exactly one merger from state $n$ is
\[
 P_{n,n-1}(v)
 =\int_0^v\lambda_ne^{-\lambda_n s}
                           e^{-\lambda_{n-1}(v-s)}\,ds
 =\frac{\lambda_n}{n-1}
       \left(e^{-\lambda_{n-1}v}-e^{-\lambda_nv}\right),
\]
because $\lambda_n-\lambda_{n-1}=n-1$.  Tonelli's theorem and the Markov
property give
\[
 R_n(f)=\int_0^L f(u)e^{-\lambda_nu}P_{n,n-1}(L-u)\,du,
\]
which is exactly \eqref{restore:Rformula}; rearrangement gives
\eqref{restore:Fformula}.  In particular $R_2=F(1)-e^{-L}T$, confirming
the normalization at the first index.

The two-lineage no-merger event has probability $e^{-\lambda_2L}=e^{-L}$,
so it determines $L$.  Let $\alpha=e^{-L}$ and define the probability
measure
\begin{equation}\label{restore:clockmeasure}
 \mu_f=(u\mapsto e^{-u})_\#\left(\frac{f(u)}{T}\,du\right)
               \quad\hbox{on }[\alpha,1].
\end{equation}
Its moments are $\int x^j\,d\mu_f=F(j)/T$ for all $j\ge0$, with
$F(0)=T$.  Formula \eqref{restore:Fformula} therefore determines all
its polynomial moments.  Polynomials are uniformly dense in
$C([\alpha,1])$, so the finite measure $\mu_f$ is unique.  The inverse
map $u=-\log x$ identifies $f(u)\,du$; continuity then identifies $f$
pointwise.  Finally $t=a+\int_0^u f(v)\,dv$ reconstructs the calendar
clock and its inverse on the interval.
\end{proof}

Thus the relevant spectral distinction is
\[
 \{\lambda_n:n\ge1\}\quad\hbox{versus}\quad
 \{\lambda_n-\lambda_{n-1}:n\ge2\}=\{1,2,3,\ldots\}.
\]
The first set admits the compactly supported annihilators used above.
The second yields a determining polynomial moment sequence after
exponentiation.  No claim of equality or recovery of a full genealogy law
is needed for this conclusion.

\subsection{The exact finite-correlation order}

The identification theorem uses arbitrarily many initial lineage counts.
At a finite maximum count, a matching upper and lower bound quantifies
the remaining ambiguity.  All Wasserstein distances below use the
ordinary Euclidean distance in the indicated coordinate.

Let $\mathcal F_{L,T}$ be the positive smooth functions on $[0,L]$ with
integral $T$.  For an integer $m\ge1$, define the worst-fiber diameter
\[
 D_m^x(L,T)=\sup\left\{W_1(\mu_f,\mu_g):
   f,g\in\mathcal F_{L,T},\quad R_n(f)=R_n(g)\quad(2\le n\le m+1)
                                                       \right\}.
\]
This supremum ranges over all possible common data vectors.  A lower
bound for it does not assert the same ambiguity in every individual
data fiber.

\begin{theorem}[Sharp finite-correlation ambiguity]
\label{restore:finite}
There are numerical constants $c,C>0$ such that, for every $L,T>0$ and
every integer $m\ge1$,
\begin{equation}\label{restore:diameter}
 c\frac{1-e^{-L}}{m+1}\le D_m^x(L,T)
                \le C\frac{1-e^{-L}}{m+1}.
\end{equation}
The upper bound holds in every data fiber.  The lower bound can be
realized by positive smooth clock densities that agree on neighborhoods
of both endpoints of $[0,L]$.

For $\nu_f=f(u)\,du/T$ on $[0,L]$ one has
\begin{equation}\label{restore:metriccompare}
 W_1(\mu_f,\mu_g)\le W_1(\nu_f,\nu_g)
                              \le e^L W_1(\mu_f,\mu_g).
\end{equation}
In particular the worst-fiber diameter in coalescent time has order
$1/(m+1)$ for every fixed $L>0$.
\end{theorem}
\begin{proof}
By \eqref{restore:Fformula}, equal correlations through $R_{m+1}$ are
equivalent to equal moments of $\mu_f$ and $\mu_g$ through degree $m$.
Write $d=1-\alpha$, where $\alpha=e^{-L}$.  The classical Jackson
approximation theorem \cite{DL}, rescaled from $[0,1]$
to $[\alpha,1]$, gives, for every 1-Lipschitz function $\varphi$, a
polynomial $P$ of degree at most $m$ satisfying
\[
                 \|\varphi-P\|_\infty\le C_Jd/m.
\]
Its integrals under the two measures agree.  Thus
$|\int\varphi\,d(\mu_f-\mu_g)|\le2C_Jd/m$.  The
Kantorovich--Rubinstein dual formula gives the upper bound in
\eqref{restore:diameter}, after adjusting a numerical constant.

For completeness the matching lower bound can be made inside the smooth
positive class.  Put $r=m+1$ and
\[
 x_j=\alpha+\frac d4+\frac{jd}{2r},\qquad 0\le j\le r,
 \qquad h=\frac d{2r}.
\]
Let
\[
 P_r=2^{1-r}\sum_{\substack{0\le j\le r\\j\ {\mathrm{even}}}}
                        \binom rj\delta_{x_j},\qquad
 Q_r=2^{1-r}\sum_{\substack{0\le j\le r\\j\ {\mathrm{odd}}}}
                        \binom rj\delta_{x_j}.
\]
Both have mass one.  The $r$th finite difference of any polynomial of
degree at most $r-1=m$ is zero, so their moments through degree $m$
coincide.  Moreover
\[
 \sum_{j=0}^i(-1)^j\binom rj=(-1)^i\binom{r-1}{i}
                           \quad(0\le i<r)
\]
and the one-dimensional CDF formula imply
\begin{equation}\label{restore:atomicdistance}
                      W_1(P_r,Q_r)=h.
\end{equation}

Choose a nonnegative smooth probability density $\kappa_\delta$ supported
in $[-\delta,\delta]$, with $0<\delta\le h/8$.  Convolution preserves
equality of polynomial moments through degree $m$: integrating a
polynomial at $x+z$ gives, for each fixed $z$, another polynomial of
the same degree in $x$.  Both convolved measures remain supported in
the interior of $[\alpha,1]$.  Coupling each atom with its translated
smooth bump shows
\[
 W_1(P_r,P_r*\kappa_\delta)\le\delta,
 \qquad W_1(Q_r,Q_r*\kappa_\delta)\le\delta.
\]
By \eqref{restore:atomicdistance}, their convolved distance is at least
$h-2\delta\ge3h/4$.

Let $U$ be the uniform probability measure on $[\alpha,1]$, and set
\[
 \mu=\tfrac12U+\tfrac12(P_r*\kappa_\delta),\qquad
 \widetilde\mu=\tfrac12U+\tfrac12(Q_r*\kappa_\delta).
\]
These measures have smooth densities $p,\widetilde p\ge1/(2d)$ on
the closed interval.  They agree near the endpoints and still have the
same first $m$ moments.  Their CDF difference is one half of the
preceding difference, so
\[
 W_1(\mu,\widetilde\mu)\ge\frac{3h}{8}
                                  =\frac{3d}{16(m+1)}.
\]
The functions
\[
 f(u)=Te^{-u}p(e^{-u}),\qquad
 g(u)=Te^{-u}\widetilde p(e^{-u})
\]
belong to $\mathcal F_{L,T}$, agree near $0,L$, and push forward to
$\mu,\widetilde\mu$ as in \eqref{restore:clockmeasure}.  They give the
lower bound.  Finally $u\mapsto e^{-u}$ is 1-Lipschitz on $[0,L]$ and
$x\mapsto-\log x$ is $e^L$-Lipschitz on $[e^{-L},1]$, proving
\eqref{restore:metriccompare}.
\end{proof}

\begin{corollary}[An exact calendar-clock interpretation]
\label{restore:clockLone}
For $f,g\in\mathcal F_{L,T}$ put
$\tau_f(u)=\int_0^u f(v)\,dv$ and
$\tau_g(u)=\int_0^u g(v)\,dv$.  Then
\[
     \|\tau_f-\tau_g\|_{L^1(0,L)}=T\,W_1(\nu_f,\nu_g).
\]
Thus the finite-correlation diameter controls the integrated discrepancy
between inverse calendar clocks, with matching order $1/(m+1)$ for fixed
$L,T$.  This is not a bound for population-size functions in a
pointwise or derivative norm.
\end{corollary}
\begin{proof}
The distribution functions of $\nu_f$ and $\nu_g$ are $\tau_f/T$ and
$\tau_g/T$.  Apply the one-dimensional CDF formula for $W_1$.
\end{proof}

\subsection{Exact identification and numerical conditioning}

If $T,L$ are known exactly and two correlation vectors have raw errors
$\Delta R_{j+1}$, their induced moment errors satisfy the exact identity
\begin{equation}\label{restore:noise}
 \left|\Delta\left(\frac{F(j)}T\right)\right|
       =\frac{2}{T(j+1)}e^{j(j-1)L/2}|\Delta R_{j+1}|.
\end{equation}
The factor is large at high lineage counts.  The identification theorem
therefore does not by itself establish a favorable finite-precision or
statistical reconstruction rate.  Conversely, this amplification formula
does not prove the absence of every uniform continuity modulus for the
inverse problem.  Theorem~\ref{restore:noisy} below gives a matching
upper and lower deterministic precision law for these correlations.
Errors in the separately estimated $L$ would also enter
\eqref{restore:Fformula} and are not included in
\eqref{restore:noise}.



\subsection{A sharp deterministic precision law for the restored clock}
\label{restore:noisysection}

For fixed, exactly known $L,T>0$, normalize the correlations by
$r_n=R_n/T$.  Suppose that the first $M$ adjacent correlations are
observed with bounded errors,
\[
       y_j=r_{j+1}(f)+e_j,\qquad |e_j|\le\varepsilon,
                              \qquad 1\le j\le M.
\]
Here $M$ counts distinct initial-lineage experiments; it is not the
number of random genealogies observed.  Let $\mathcal R_{M,L}(\varepsilon)$
be the smallest possible worst-case error in $W_1$ for recovery of
$\mu_f$, with supremum over $f\in\mathcal F_{L,T}$ and all the displayed
error sequences.  An estimator may return any probability measure on
$[e^{-L},1]$.  The case $M=\infty$ means that all normalized correlations
are supplied with the same coordinatewise error bound.

Equivalently, raw observations of $R_{j+1}$ with absolute errors at most
$\delta$ give this model with $\varepsilon=\delta/T$.  This normalization
keeps the logarithm dimensionless and separates the known calendar
duration from the precision parameter.

\begin{theorem}[Finite-correlation and finite-precision resolution]
\label{restore:noisy}
For every fixed $L>0$ there exist constants $c_L,C_L>0$ such that,
for $0<\varepsilon\le e^{-2}$ and $M\in\mathbb N\cup\{\infty\}$,
\[
 \frac{c_L}{1+\min\{M,\sqrt{\log(1/\varepsilon)}\}}
 \ \le\ \mathcal R_{M,L}(\varepsilon)
 \ \le\
 \frac{C_L}{1+\min\{M,\sqrt{\log(1/\varepsilon)}\}}.
\]
In particular, supplying every adjacent correlation does not improve
the precision-limited order $1/\sqrt{\log(1/\varepsilon)}$ in this
deterministic observation model.  Constants are allowed to depend on
$L$; this statement is not a uniform small-gap or large-gap theorem.
For raw errors $0<\delta\le Te^{-2}$, the same conclusion holds with
$\log(1/\varepsilon)$ replaced by $\log(T/\delta)$.
\end{theorem}
\begin{proof}
Put $\alpha=e^{-L}$ and $d=1-\alpha$.  A coefficient-controlled form
of Jackson approximation gives, for any two probability measures
$\mu,\nu$ on $[\alpha,1]$ and any integer $k\ge1$,
\begin{equation}\label{restore:momentstable}
 W_1(\mu,\nu)\le\frac{A_L}{k}
       +A_L B_L^k\max_{1\le j\le k}
                         \left|\int x^j\,d(\mu-\nu)\right|,
\end{equation}
with finite constants $A_L,B_L>1$.  To verify the coefficient control,
subtract the value of a 1-Lipschitz test function at $\alpha$ and apply
Jackson approximation on $[\alpha,1]$.  The approximating polynomial
can be chosen with norm at most a numerical multiple of $d$.  Expand
it in the shifted Chebyshev basis
$T_j((2x-1-\alpha)/d)$; its Chebyshev coefficients are bounded by a
numerical multiple of $d$.  The three-term recursion shows that the
sum of absolute monomial coefficients of the $j$th shifted Chebyshev
polynomial is at most $(9/d)^j$.  Summing the coefficients and using
the dual formula for $W_1$ proves \eqref{restore:momentstable}.

Choose any measure whose first $k\le M$ predicted normalized
correlations fit the data with residual at most $\varepsilon$.
The true measure is feasible.  By \eqref{restore:noise}, its moments
and those of the fit satisfy
\[
 \max_{1\le j\le k}|\Delta m_j|
           \le 2\varepsilon e^{Lk(k-1)/2}.
\]
Thus the error is bounded by
\[
 \frac{A_L}{k}
               +2A_L\varepsilon B_L^ke^{Lk(k-1)/2}.
\]
With $H=\log(1/\varepsilon)$, choose
$k=\lfloor c\min\{M,\sqrt H\}\rfloor$ for a sufficiently small
constant $c>0$ depending only on $L,A_L,B_L$.  When $k\ge1$, the
exponential term is bounded by a constant times $e^{-H/2}$ and the
first term has the claimed order.  When this choice gives $k=0$,
the trivial bound $W_1\le d$ supplies the same conclusion after
enlarging $C_L$.

For a lower bound, exact matching of the first $M$ correlations in
Theorem~\ref{restore:finite} gives error at least a constant times
$d/(M+1)$; omit this bound when $M=\infty$.  For the precision
obstruction, use the smooth positive pair constructed there with
moments matching through an integer $m$.  Its $W_1$ distance is at
least $3d/(16(m+1))$.  Its normalized correlations coincide for
$n\le m+1$.  For $n\ge m+2$, since both moment values lie in $[0,1]$,
\[
 |r_n(f)-r_n(g)|
   =\frac n2 e^{-\lambda_{n-1}L}
       \left|\int x^{n-1}\,d(\mu_f-\mu_g)\right|
   \le n\exp\left(-\frac{L(n-1)(n-2)}2\right).
\]
There is a number $N_L$ such that
$\log n\le L(n-2)^2/4$ for every $n\ge N_L$.  Choose
\[
          m=\left\lceil N_L+2\sqrt{H/L}\right\rceil.
\]
The displayed correlation difference is then at most $e^{-H}$ for
every $n\ge m+2$, and hence at most $2\varepsilon$ for every
$n\ge2$.  The coordinatewise midpoint of the two complete correlation
vectors is consistent with either clock at the permitted error level.
Any estimator must therefore incur at least half their $W_1$ distance,
which is at least $c_L/(1+\sqrt H)$.  Combining this with the exact
finite-correlation lower bound proves the theorem.
\end{proof}

This result closes the deterministic noise question for these specified
correlations with exactly known $L,T$.  It does not give the number of
random loci needed to estimate them, account for uncertainty in $L$,
or assert that endpoint--occupation correlations can be read from an
ordinary site-frequency spectrum.


\section{A positive boundary: completely monotone histories}
\label{sec:complete-monotonicity}

The failure of identifiability on a smooth unrestricted class does not
extend to every infinite-dimensional shape class.  We give a sufficient
condition that uses all derivative signs simultaneously.  This condition
differs from imposing any fixed finite list of strict differential
inequalities, which a sufficiently small localized perturbation can preserve.

Write
\[
 C_\eta(t)=\int_0^t\frac{ds}{\eta(s)},\qquad
 T_\eta=C_\eta^{-1},\qquad
 f_\eta(u)=T_\eta'(u)=\eta(T_\eta(u)).
\]
A function $f$ is completely monotone if
$(-1)^j f^{(j)}(u)\geq0$ for every integer $j\geq0$ and $u>0$.
By the Bernstein representation theorem, a bounded completely monotone
function continuous at zero has the representation
\begin{equation}\label{eq:cm-bernstein}
 f(u)=\int_{[0,\infty)} e^{-ru}\,d\mu(r)
\end{equation}
for a unique finite positive measure $\mu$; see
\cite[Chapter~1]{SSV2012}.  The additional conditions
\begin{equation}\label{eq:cm-measure-class}
 \mu(\{0\})>0,\qquad
 \int r^j\,d\mu(r)<\infty\quad(j=1,2,\ldots)
\end{equation}
ensure that $f\in C^\infty([0,\infty))$ and
$0<\mu(\{0\})\leq f\leq\mu([0,\infty))$.
Consequently $T(u)=\int_0^u f(v)\,dv$ is a smooth increasing bijection
of $[0,\infty)$, and
\begin{equation}\label{eq:cm-history-reconstruction}
 \eta(t)=f(T^{-1}(t))
\end{equation}
is a smooth positive population history bounded above and away from zero.

\begin{lemma}[A zero-set criterion in a half-plane]\label{lem:cm-zeros}
Let $H$ be bounded and holomorphic on $\{w:\operatorname{Re}w>0\}$.
Suppose $H(a_k)=0$ for distinct positive numbers $a_k\to\infty$ and
$\sum_k a_k^{-1}=\infty$.  Then $H$ vanishes identically.
\end{lemma}
\begin{proof}
For $a>0$ put $b_a(w)=(w-a)/(w+a)$.
If $H(a)=0$, the Schwarz lemma after mapping the right half-plane to
the unit disk shows that $H/b_a$ is holomorphic and has supremum norm
at most $\|H\|_\infty$.  Repeated application therefore gives, for every
$N$ and every positive $w$ distinct from the $a_k$,
\[
 |H(w)|\leq \|H\|_\infty
       \prod_{k=1}^N\left|\frac{w-a_k}{w+a_k}\right|.
\]
For $a_k>2w$,
\[
 \log\frac{a_k-w}{a_k+w}
 =\log(1-w/a_k)-\log(1+w/a_k)\leq-2w/a_k.
\]
The product tends to zero.  Thus $H$ vanishes on the positive real
axis, and the identity theorem finishes the proof.
\end{proof}

\begin{theorem}[Uniqueness of positive Stieltjes mixtures]
\label{thm:cm-stieltjes}
Let $(\rho_k)$ be a sequence of distinct positive numbers tending to
infinity and satisfying
\begin{equation}\label{eq:cm-spectrum-condition}
 \sum_k\rho_k^{-1/2}=\infty.
\end{equation}
If finite positive measures $\mu_1,\mu_2$ on $[0,\infty)$ satisfy
\[
 \int\frac{d\mu_1(r)}{\rho_k+r}
 =\int\frac{d\mu_2(r)}{\rho_k+r}
 \qquad\text{for every }k,
\]
then $\mu_1=\mu_2$.  Condition~\eqref{eq:cm-spectrum-condition} is a
sufficient condition; no necessity assertion is made here.
\end{theorem}
\begin{proof}
For a finite positive measure $\mu$ define
\[
 F_\mu(z)=\int\frac{d\mu(r)}{z+r},\qquad
 G_\mu(w)=wF_\mu(w^2)
 =\frac12\int\left(\frac1{w-i\sqrt r}
                    +\frac1{w+i\sqrt r}\right)d\mu(r).
\]
The last expression is holomorphic for $\operatorname{Re}w>0$:
on each compact subset the integrand and its complex derivatives
are bounded uniformly in $r$.  Both summands have positive real part.
It follows that $\operatorname{Re}G_\mu\geq0$, and
\[
 B_\mu(w)=\frac{G_\mu(w)-1}{G_\mu(w)+1}
\]
is holomorphic and bounded by one in modulus.  This also covers the
zero measure, for which $B_\mu\equiv-1$.
The matching Stieltjes values imply
$B_{\mu_1}(\sqrt{\rho_k})=B_{\mu_2}(\sqrt{\rho_k})$.
Lemma~\ref{lem:cm-zeros} applied to their difference yields
$B_{\mu_1}=B_{\mu_2}$, hence $G_{\mu_1}=G_{\mu_2}$ and
$F_{\mu_1}(z)=F_{\mu_2}(z)$ for every positive $z$.

For completeness, this last equality determines the measure without
any moment assumptions at infinity.  Put $\nu=\mu_1-\mu_2$.
Dominated convergence applied to $z/(z+r)$ as $z\to\infty$ gives
$\int d\nu=0$.  Differentiating the zero Stieltjes transform at $z=1$
gives
\[
 \int (1+r)^{-m}\,d\nu(r)=0\qquad(m=1,2,\ldots).
\]
Push $\nu$ forward by $r\mapsto(1+r)^{-1}$ to a finite signed measure
on $[0,1]$, giving it zero mass at zero.  Every polynomial integral
vanishes.  Density of polynomials in $C([0,1])$ implies that this
measure, and therefore $\nu$, is zero.
\end{proof}

\begin{theorem}[Identification from one sampling date]
\label{thm:cm-coalescent}
Consider the class of histories for which the inverse-clock derivative
$f_\eta$ has representation~\eqref{eq:cm-bernstein} with
\eqref{eq:cm-measure-class}.  Within this class the absolute expected
unfolded SFS at sampling date zero, with known constant mutation rate
and known for every sample size
$n\geq2$, identifies $\eta$ uniquely.  Equivalently, the first-coalescence
means $c_k$ for all $k\geq2$ identify $\eta$ uniquely.
\end{theorem}
\begin{proof}
Put $\lambda_k=\binom{k}{2}$.  The survival probability of the first
coalescence among $k$ contemporaneous lineages is
$\exp(-\lambda_k C_\eta(t))$, so its mean is
\begin{align*}
 c_k
 &=\int_0^\infty \exp(-\lambda_k C_\eta(t))\,dt
   =\int_0^\infty e^{-\lambda_k u} f_\eta(u)\,du\\
 &=\int_{[0,\infty)}\frac{d\mu(r)}{\lambda_k+r}.
\end{align*}
The last step follows from Tonelli's theorem.
Since $\sqrt{\lambda_k}\sim k/\sqrt2$, the reciprocal square-root
sum diverges.  Theorem~\ref{thm:cm-stieltjes} therefore identifies
$\mu$, hence $f_\eta$, then $T_\eta(u)=\int_0^u f_\eta(v)\,dv$,
and finally $\eta$ by~\eqref{eq:cm-history-reconstruction}.

To connect the two stated data sets, Bhaskar and Song
\cite[Lemma~1 and its proof]{BS} establish that the absolute
expected branch-length SFS for sample size $n$ determines
$(c_2,\ldots,c_n)$ by an invertible linear transformation.  Their
lemma is stated for a normalized spectrum, hence up to one common
scale; their proof supplies the absolute version used here.
A known mutation rate converts absolute expected mutation counts
into the absolute branch-length spectrum.
\end{proof}

\begin{corollary}[Complete monotonicity in calendar time]
\label{cor:cm-calendar}
The conclusion of Theorem~\ref{thm:cm-coalescent} holds for the class
of histories $\eta\in C^\infty([0,\infty))$ that are bounded above
and away from zero and satisfy
$(-1)^j\eta^{(j)}(t)\geq0$ for all $j\geq0$ and $t\geq0$.
\end{corollary}
\begin{proof}
We verify that complete monotonicity of $\eta$ implies complete
monotonicity of $f_\eta$.  Since $T_\eta'=\eta\circ T_\eta$,
\[
 f_\eta^{(n)}(u)
 =\left[(\eta(t)D_t)^n\eta(t)\right]_{t=T_\eta(u)}.
\]
The differential polynomial $(\eta D_t)^n\eta$ is a sum of monomials
with nonnegative coefficients, and in each monomial the sum of the
orders of all differentiated factors is $n$.  This follows by induction:
at the next step one differentiates one factor and multiplies by
$\eta$.  Every monomial consequently has sign $(-1)^n$ or is zero.
Thus $f_\eta$ is completely monotone.

The Bernstein representation has finite mass $f_\eta(0)$.
Its atom at zero equals $\lim_{u\to\infty}f_\eta(u)>0$.
For $j\geq1$, differentiation for $u>0$ and monotone convergence as
$u\downarrow0$ give
$\int r^j\,d\mu(r)=(-1)^j f_\eta^{(j)}(0)<\infty$.
Therefore~\eqref{eq:cm-measure-class} holds, and
Theorem~\ref{thm:cm-coalescent} applies.
\end{proof}

\begin{remark}[Scope of the positive result]
This is an exact uniqueness result for an infinite collection of
expectations, not a finite-sample or stability guarantee.  Complete
monotonicity imposes infinitely many weak inequalities and should not
be confused with imposing monotonicity, convexity, or any other fixed
finite collection of derivative signs.  Nor does the conclusion hold
against arbitrary competitors outside the stated shape class.
The half-plane argument is a classical uniqueness mechanism for
positive Stieltjes transforms.  We use it to establish the stated
shape-restricted boundary and do not claim a new complex-analytic
uniqueness theorem.
\end{remark}



\section{Joint recovery of demography and mutation intensity}
\label{sec:mutation}

The temporal-design theorem was stated with a known constant mutation
rate so that expected spectra could be interpreted immediately as
branch lengths. Heterochronous pairs supply a direct calibration that
removes this requirement. Throughout this section the expected joint
spectrum is unfolded: the ancestral state is distinguished, and the
two sampling cohorts remain labeled.

\begin{proposition}[Calibration by a heterochronous pair]
\label{prop:mutationconstant}
Take one lineage at date $a$ and one at date $b>a$. Let $Y_a,Y_b$
denote the counts of mutations private to the respective sampled
lineages. If the mutation rate per unit calendar branch length is
an unknown constant $\mu>0$, then
\[
 \E Y_a-\E Y_b=\mu(b-a).
\]
Consequently any two distinct allowed dates identify $\mu$ from
their absolute expected joint spectrum. Theorem~\ref{phy:design}
therefore also gives an injectivity criterion for the pair
$(\eta,\mu)$: it is identifiable on
$\mathcal H\times(0,\infty)$ if and only if $S$ is dense.
\end{proposition}
\begin{proof}
If the common ancestor occurs at calendar time $\tau\ge b$, the two
private branches have lengths $\tau-a$ and $\tau-b$. Their difference
is exactly $b-a$, before expectation. The infinite-sites expectation
is mutation rate times branch length, proving the formula. A dense
set contains two distinct dates, so after calibration the positive
part of Theorem~\ref{phy:design} applies. A nondense set remains
noninjective even if the same constant $\mu$ is fixed under both
histories.
\end{proof}

\begin{theorem}[An unknown time-dependent mutation clock]
\label{thm:mutationvariable}
Let both the population history $\eta$ and the mutation intensity
$\rho$ be smooth functions with positive uniform lower and upper
bounds on $[0,\infty)$. Suppose $\rho(t)$ is the same deterministic
calendar-time intensity on every extant lineage. The collection of
all absolute expected joint serial spectra at dates in $S$ identifies
$(\eta,\rho)$ if and only if $S$ is dense.
\end{theorem}
\begin{proof}
For one sample at $a$ and one at $b>a$, the preceding pathwise
argument gives
\begin{equation}\label{eq:mutationdifference}
 D(a,b):=\E Y_a-\E Y_b=\int_a^b\rho(v)\,dv.
\end{equation}
Knowledge of these integrals for $a,b$ in a dense set determines the
continuous function $\rho$. Indeed continuity extends the integrals
to every finite interval, and differentiation recovers its integrand.

For a contemporaneous pair at date $t$, the expectation of private
mutations on either branch is
\[
 g(t)=\int_t^\infty\rho(v)
       \exp\left(-\int_t^v\frac{dw}{\eta(w)}\right)dv.
\]
Boundedness and positivity justify differentiation, giving
\begin{equation}\label{eq:jointinverse}
 g'(t)=\frac{g(t)}{\eta(t)}-\rho(t),\qquad
 \eta(t)=\frac{g(t)}{g'(t)+\rho(t)}.
\end{equation}
Dense contemporaneous dates determine the continuous function $g$,
and hence its derivative, so they determine $\eta$ after $\rho$ has
been recovered from \eqref{eq:mutationdifference}. The denominator
in \eqref{eq:jointinverse} is positive because it equals $g/\eta$.
For necessity, restrict to the subclass with $\rho\equiv1$ and use
the gap family of Theorem~\ref{phy:design}.
\end{proof}

The time-dependent assertion assumes an intensity shared by all
lineages and known calendar dates. It does not address lineage-specific
mutation rates, selection, population structure, ascertainment, or
unknown temporal calibration. The negative argument for variable
$\rho$ is a counterexample in its constant-intensity subclass; it is
not a claim that the original gap perturbation preserves an arbitrary
fixed nonconstant mutation intensity.

\begin{proposition}[A joint clock ambiguity inside a sampling gap]
\label{prop:jointgauge}
Let $I$ be a sampling gap, and let $\phi$ be a smooth increasing
diffeomorphism of $[0,\infty)$ equal to the identity outside a compact
subinterval of $I$. For a baseline pair $(\eta_0,\rho_0)$ as in
Theorem~\ref{thm:mutationvariable}, define
\[
 \eta_\phi(t)=\frac{\eta_0(\phi(t))}{\phi'(t)},\qquad
 \rho_\phi(t)=\rho_0(\phi(t))\phi'(t).
\]
These pairs have identical laws of the mutation-marked genealogical
tree after its calendar event times and calendar edge lengths have been
forgotten, for every
finite sampling schedule on $S$. In particular they have identical
single-locus infinite-sites observation laws. Nontrivial $\phi$
arbitrarily close to the identity give distinct smooth pairs.
\end{proposition}
\begin{proof}
Under $s=\phi(t)$ the coalescent rate and mutation intensity satisfy
\[
 \frac{dt}{\eta_\phi(t)}=\frac{ds}{\eta_0(s)},\qquad
 \rho_\phi(t)\,dt=\rho_0(s)\,ds.
\]
Thus one may couple every coalescent merger and every mutation mark
by this deterministic time change. All sampling dates are fixed by
$\phi$, so the insertions, descendant labels, merger topology, and
mutation assignments agree. Positivity and boundedness persist because
$\phi'$ is positive, bounded, and equal to one off a compact set.
If the population history alone were unchanged, its integrated
coalescent clock would satisfy $C_0(\phi(t))=C_0(t)$; strict
monotonicity would force $\phi(t)=t$. Hence a nontrivial time change
gives a distinct pair.
\end{proof}

This is the usual deterministic time-change symmetry applied jointly
to the two clocks. Its equality of full mutation observations is
stronger than the fixed-mutation expectation result, because the
mutation intensity is now allowed to vary with the demographic
history. It must not be transferred to a model with a fixed known
mutation clock or a fixed calendar-time recombination process.



\section{What is identified, and what remains to be determined}
\label{sec:discussion}

The results distinguish three observation mechanisms. Ordinary expected
branch rewards see the coalescent exit rates. Endpoint-conditioned
occupation rewards see differences of exit rates. Heterochronous
private-mutation differences isolate the calendar mutation clock.
These mechanisms are mathematically distinct, even when they involve
the same underlying genealogy.

The negative theorem is not a claim that ancient samples are useless.
It says that unrestricted smooth demography cannot be identified
\emph{inside every sampling gap} from the entire collection of specified
first expectations, without additional assumptions. A fixed-complexity
parametric family, complete monotonicity, richer genealogical
observations, or an observation process that itself depends on demography
can change the conclusion. Conversely, pointwise identifiability under
dense sampling does not furnish a stable numerical estimator from a
finite genome.

The finite-correlation and noisy-correlation theorems address a
deterministic inverse problem with known interval length in both clocks.
The observed quantities are scalar correlations, not raw genealogical
samples. Converting their precision requirement into a minimax rate for
a specified genomic experiment requires control of their estimation
variance and dependence. The population-size function is a derivative
of the inverse clock; Wasserstein control of the clock measure alone
does not imply uniform control of that derivative.

Several further questions remain concrete. One is a sharp joint
stability theory when the coalescent interval length $L$ is also measured
with error. Another is to determine which finite-locus spectrum
covariances expose a complete spectral set without access to the
occupation correlations used here. For general coalescents, the
annihilator method requires an actual spectral expansion, control of
Jordan multiplicities, and a separate examination of which spectral
differences are observable. Preferential sampling, recombination,
population structure, and selection require their own models; the
present proofs do not resolve those extensions.



\appendix


\section{Compact nullspaces for sparse spectral observations}
\label{spec:section}\label{sec:sparse}

For a compactly supported integrable function $h$ on $\mathbb R$, write
\[
 \mathcal Lh(z)=\int_{\mathbb R}e^{-zu}h(u)\,du,
 \qquad z\in\mathbb C.
\]
The same notation is used for a finite signed or complex measure of
compact support. Its transform is entire. A zero of multiplicity $m$
at $z=\lambda$ is equivalent to the $m$ conditions
\[
 \int_{\mathbb R}u^j e^{-\lambda u}h(u)\,du=0,
 \qquad 0\le j<m.
\]
These are precisely the exponential--polynomial terms generated by
Jordan blocks of length at most $m$.

\begin{theorem}[Local annihilation for a superlinear spectrum]
\label{spec:sparse}
Let $\alpha>1$, $c>0$, and let $(\lambda_k)_{k\ge1}$ be positive real
numbers satisfying $\lambda_k\ge c k^\alpha$. Fix an integer $m\ge1$,
a bounded nonempty open interval $J\subset\mathbb R$, and a real number
$\sigma$ with $1<\sigma<\alpha$. There is an infinite-dimensional
real vector space of functions $h\in C_c^\infty(J)$, all of Gevrey
order $\sigma$, such that
\begin{align}
 \int u^j h(u)\,du&=0 &&(0\le j<m),\label{spec:zero-moments}\\
 \int u^j e^{-\lambda_k u}h(u)\,du&=0
       &&(k\ge1,\ 0\le j<m).\label{spec:exponential-moments}
\end{align}
Intersecting this space with the kernels of any finite collection of
real linear functionals on $C_c^\infty(J)$ still leaves an
infinite-dimensional space. In particular, any finite collection of
additional spectral-jet conditions can be imposed simultaneously.
\end{theorem}

\begin{proof}
Define
\begin{equation}\label{spec:product}
 F(z)=z^m\prod_{k=1}^{\infty}
                 \left(1-\frac z{\lambda_k}\right)^m.
\end{equation}
The product converges uniformly on compact subsets of $\mathbb C$,
because $\sum_k\lambda_k^{-1}<\infty$. It defines a nonzero real
entire function, with a zero of order at least $m$ at every
$\lambda_k$ and at zero.

For $R\ge0$, monotonicity of the summand gives
\begin{align*}
 \sum_{k=1}^{\infty}\log\left(1+\frac R{\lambda_k}\right)
 &\le\sum_{k=1}^{\infty}\log\left(1+\frac R{c k^\alpha}\right)\\
 &\le\int_0^\infty\log\left(1+\frac R{c x^\alpha}\right)\,dx\\
 &=c^{-1/\alpha}R^{1/\alpha}
          \int_0^\infty\log(1+y^{-\alpha})\,dy.
\end{align*}
The last integral is finite: its integrand is logarithmic at zero
and is $O(y^{-\alpha})$ at infinity. Absorbing the finite factor
$R^m$ into the exponential, we obtain constants $C_0,C_1>0$ such that
\begin{equation}\label{spec:growth}
                   |F(z)|\le C_0\exp(C_1|z|^{1/\alpha}).
\end{equation}
Write $F(z)=\sum_{j\ge0}a_jz^j$. Cauchy's coefficient estimate on
the circle of radius $R=(\alpha j/C_1)^\alpha$, for $j\ge1$, yields
\[
 |a_j|\le C_0
       \left(\frac{eC_1}{\alpha j}\right)^{\alpha j}
       \le \frac{C_0 B^j}{(j!)^\alpha}
\]
for a fixed $B>0$, since $j!\le j^j$. Enlarging $C_0$ includes $j=0$.

Choose a nonnegative, nonzero Gevrey-$\sigma$ bump $\kappa$ whose
support is a compact subinterval of $J$. Thus, for some $M,A>0$,
\[
                    \|\kappa^{(j)}\|_\infty
                         \le M A^j(j!)^\sigma\qquad(j\ge0).
\]
Consider the infinite-order differential expression
\begin{equation}\label{spec:differential}
                  h=F(D)\kappa=\sum_{j\ge0}a_j\kappa^{(j)},
                  \qquad D=\frac d{du}.
\end{equation}
For every integer $q\ge0$, the bound
$(j+q)!\le2^{j+q}j!q!$ implies
\[
 |a_j|\,\|\kappa^{(j+q)}\|_\infty
 \le C_0M(2^\sigma A)^q(q!)^\sigma
                \frac{(2^\sigma AB)^j}{(j!)^{\alpha-\sigma}}.
\]
The last series is summable and its sum is independent of $q$.
Consequently \eqref{spec:differential} converges together with every
derivative, defines a Gevrey-$\sigma$ function, and has support
contained in $\operatorname{supp}\kappa$.

Uniform convergence on this common compact support permits
termwise integration. Repeated integration by parts has no boundary
terms and gives, for every $z\in\mathbb C$,
\begin{equation}\label{spec:multiplier}
                         \mathcal Lh(z)=F(z)\mathcal L\kappa(z).
\end{equation}
The prescribed zeros and their multiplicities prove
\eqref{spec:zero-moments}--\eqref{spec:exponential-moments}.
Moreover $F(-1)\ne0$ and $\mathcal L\kappa(-1)>0$, so $h\ne0$.

Every derivative $D^qh$ has the same support inclusion and satisfies
the same annihilation conditions, because its Laplace transform is
$z^q\mathcal Lh(z)$. These derivatives are linearly independent:
a nontrivial finite linear relation would give
$Q(z)\mathcal Lh(z)=0$ for a nonzero polynomial $Q$, contradicting
the nonzero entire function in \eqref{spec:multiplier}.
Their real linear span is therefore an infinite-dimensional space
with the required properties. A linear map from this span to
$\mathbb R^d$ has a kernel of codimension at most $d$; applying this
to any finite collection of functionals proves the last assertion.
\end{proof}

\begin{remark}[Compact Gevrey bumps]
\label{spec:bump}
The bump used above can be constructed directly. Put
$p=(\sigma-1)^{-1}$ and define $f(x)=\exp(-x^{-p})$ for $x>0$,
with $f(x)=0$ for $x\le0$. On a complex circle
$|z-x|=\delta x$, where $\delta>0$ depends only on $p$, the branch
through the positive axis satisfies
$\operatorname{Re}(z^{-p})\ge c_p x^{-p}$. Cauchy's estimate gives
\[
 |f^{(j)}(x)|\le j!(\delta x)^{-j}e^{-c_p x^{-p}}.
\]
Maximizing the right-hand side over $x>0$ bounds it by
$C^j(j!)^{1+1/p}=C^j(j!)^\sigma$. The extension is flat at zero.
For $[a,b]\subset J$, the function
$\kappa(u)=f(u-a)f(b-u)$ is positive inside $(a,b)$, vanishes
outside $[a,b]$, and obeys the same Gevrey order by Leibniz's rule.
Thus the construction imposes no lower bound on the length of $J$.
\end{remark}

\begin{corollary}[Bounded Jordan multiplicities and finite calibrations]
\label{spec:jordan}
Under the hypotheses of Theorem~\ref{spec:sparse}, suppose every
observed reward on a sampling gap is a finite linear combination of
$u^je^{-\lambda_k u}$ and of $u^j$, with $0\le j<m$. Then there is
an infinite-dimensional space of smooth perturbations supported in
that gap which annihilate every such reward simultaneously. Any
finitely many additional linear calibrations can also be retained.
\end{corollary}

\begin{proof}
Apply \eqref{spec:zero-moments}--\eqref{spec:exponential-moments}
term by term, followed by the finite-codimension conclusion of
Theorem~\ref{spec:sparse}.
\end{proof}

The corollary concerns the indicated observation space. Applying it
to a particular stochastic model requires verification of that
model's spectral expansion and of the preservation of its sampling
clock. No conclusion about arbitrary unbounded Jordan lengths is
asserted.

\begin{theorem}[An obstruction for a reciprocal-divergent spectrum]
\label{spec:divergent}
Let $0<\lambda_1<\lambda_2<\cdots$ tend to infinity and suppose
\[
                         \sum_{k=1}^{\infty}\frac1{\lambda_k}
                                  =\infty.
\]
If a finite signed or complex measure $\mu$ of compact support on
$\mathbb R$ satisfies
\[
                     \int e^{-\lambda_k u}\,d\mu(u)=0
                                  \qquad(k\ge1),
\]
then $\mu=0$. In particular there is no nonzero compactly supported
smooth local annihilator for this spectrum.
\end{theorem}

\begin{proof}
Choose $a\le b$ containing the support and put
\[
 G(z)=e^{az}\mathcal L\mu(z)
     =\int_{[a,b]}e^{-z(u-a)}\,d\mu(u).
\]
For $\operatorname{Re}z\ge0$ this is bounded by the total variation
$\|\mu\|_{\mathrm{TV}}$. If $\mu\ne0$, its entire transform is
not identically zero: otherwise every polynomial moment vanishes,
and polynomial approximation on $[a,b]$ implies $\mu=0$.

Map the unit disk to the right half-plane by
$z=(1+w)/(1-w)$. The bounded analytic function
\[
                         g(w)=G\left(\frac{1+w}{1-w}\right)
\]
has zeros
$a_k=(\lambda_k-1)/(\lambda_k+1)$. Remove its finite-order zero
at $w=0$, if there is one. The quotient remains bounded on the
disk and is nonzero at zero; call it $g_0$.
Jensen's formula and boundedness give, for $0<r<1$ away from the
zero moduli,
\[
 \sum_{|a|<r}\log\frac r{|a|}
 \le \log\|g_0\|_\infty-\log|g_0(0)|,
\]
where zeros are counted with multiplicity. Letting $r\uparrow1$
shows that $\sum_a\log(1/|a|)<\infty$.
For every sufficiently large $k$, $0<a_k<1$ and
\[
 \log\frac1{a_k}\ge1-a_k=\frac2{\lambda_k+1}.
\]
Since the $\lambda_k$ are distinct, their zeros are distinct, and
the last two inequalities imply $\sum_k\lambda_k^{-1}<\infty$,
a contradiction. Thus $\mu=0$.
\end{proof}

The two theorems give a dichotomy for power-growth spectra:
superlinear growth admits local smooth nullspaces, whereas distinct
rates comparable to $k^\alpha$ with $0<\alpha\le1$ do not.
The sufficient construction in Theorem~\ref{spec:sparse} is stated
under its power-growth hypothesis; it is not a claim that the same
Gevrey construction covers every sequence with a convergent
reciprocal sum.


\begin{thebibliography}{9}
\bibitem{Kingman}
J. F. C. Kingman,
On the genealogy of large populations,
\emph{Journal of Applied Probability} \textbf{19}, suppl. A
(1982), 27--43.
\href{https://doi.org/10.2307/3213548}{doi:10.2307/3213548}.

\bibitem{MFP}
S. Myers, C. Fefferman, and N. Patterson,
Can one learn history from the allelic spectrum?,
\emph{Theoretical Population Biology} \textbf{73} (2008), 342--348.
\href{https://doi.org/10.1016/j.tpb.2008.01.001}{doi:10.1016/j.tpb.2008.01.001}.

\bibitem{BS}
A. Bhaskar and Y. S. Song,
Descartes' rule of signs and the identifiability of population
demographic models from genomic variation data,
\emph{Annals of Statistics} \textbf{42} (2014), 2469--2493.
\href{https://doi.org/10.1214/14-AOS1264}{doi:10.1214/14-AOS1264}.

\bibitem{SKS}
J. P. Spence, J. A. Kamm, and Y. S. Song,
The site frequency spectrum for general coalescents,
\emph{Genetics} \textbf{202} (2016), 1549--1561.
\href{https://doi.org/10.1534/genetics.115.184101}{doi:10.1534/genetics.115.184101}.

\bibitem{DL}
R. A. DeVore and G. G. Lorentz,
\emph{Constructive Approximation},
Grundlehren der mathematischen Wissenschaften, vol. 303,
Springer, Berlin, 1993.
\href{https://doi.org/10.1007/978-3-662-02888-9}{doi:10.1007/978-3-662-02888-9}.

\bibitem{SSV2012}
R. L. Schilling, R. Song, and Z. Vondra\v{c}ek,
\emph{Bernstein Functions: Theory and Applications}, 2nd ed.,
De Gruyter Studies in Mathematics, vol. 37,
De Gruyter, Berlin, 2012.
\href{https://doi.org/10.1515/9783110269338}{doi:10.1515/9783110269338}.
\end{thebibliography}
\end{document}
